\subsection{Разработанные улучшения}

К решению добавлена возможность выделять желтым цветом гауссианы из конкретного класса. По таблице \ref{tab:gaussian_texture} видно, что 4 байта текстуры не используются, именно их можно использовать для метки класса, не изменяя существенно оригинальный код.

Cначала требуется каждой гауссиане присвоить класс по принципу ближайшего соседа в размеченной воксельной сетке (листинг \ref{lst:segment-gaussians}). Затем к формату \texttt{.splat} добавить еще 4 байта этих классов. У оригинальной реализации был скрипт конвертера из \texttt{.ply} в \texttt{.splat}, но он медленный, так как использует циклы python, поэтому он был переписан с использованием библиотеки numpy (листинг \ref{lst:ply-to-splat}).

Изменению подверглось формирование текстуры, где теперь 12-15 байты первого пикселя содержат метку класса, а также вершинный шейдер, в который добавлена uniform переменная с выбранным классом. Каждая гауссиана выбранного класса становится желтого оттенка, а остальным уменьшается непрозрачность.

Полный код форка представлен в приложении листингами \ref{lst:splat-html} (html) и \ref{lst:splat-js} (js).